\section{Group algebras}
\label{03}

Let $K$ be a field, and $G$ be a group. The \emph{group algebra} $K[G]$ 
is the vector space (over $K$) with basis $\{g:g\in G\}$ 
and the algebra structure is given by the multiplication
\[
	\left(\sum_{g\in G}\lambda_gg\right)\left(\sum_{h\in G}\mu_hh\right)
	=\sum_{g,h\in G}\lambda_g\mu_h(gh).
\]
Every element of $K[G]$ is a finite sum of the form $\sum_{g\in G}\lambda_gg$.

\begin{exercise}
\label{xc:K[G]notsimple}
    If $G$ is non-trivial, then $K[G]$ is not simple. 
\end{exercise}

\begin{exercise}
\label{xca:K_cyclic}
  For $n\geq2$ let $G=C_n$ be the (multiplicative) cyclic group of order $n$.
  Prove that $K[G]\simeq K[X]/(X^n-1)$. 
\end{exercise}

\begin{exercise}
\label{xca:abelian=>domain}
	Let $G$ be a finitely generated torsion-free abelian group. Prove that 
	$K[G]$ is a domain. 
\end{exercise}



\begin{exercise}
	Let $G$ be a group and $\alpha=\sum_{g\in G}\lambda_gg\in K[G]$.  
	The \emph{support} of $\alpha$ is the set 
	\[
		\supp\alpha=\{g\in G:\lambda_g\ne 0\}.
	\]
	Prove that if $g\in G$, then 
	$\supp(g\alpha)=g(\supp\alpha)$ and $\supp(\alpha g)=(\supp\alpha)g$.
\end{exercise}

\begin{exercise}
\label{xca:invertible_subgroups}
	Let $G$ be a group and $H$ be a subgroup of $G$. Let $\alpha\in K[H]$. Prove that 
    $\alpha$ is invertible (resp. a left zero divisor) in $K[H]$ if and only if 
	$\alpha$ is invertible (resp. a left zero divisor) in
	$K[G]$.
\end{exercise}

\begin{exercise}
	Let $G=C_2=\langle g\rangle\simeq\Z/2$, the (multiplicative) cyclic group 
  of order two. Note that every element of $K[G]$ is of the form
	$a+bg$ for some $a,b\in K$. Prove the following statements:
	\begin{enumerate}
	    \item If the characteristic of $K$ is different from two, then 
	    \[
		K[G]\to K\times K,
		\quad
    a1_{K[G]}+bg\mapsto (a+b,a-b),
	\]
	is an algebra isomorphism. 
	\item If the characteristic of $K$ is two, then 
	\[
	K[G]\to \begin{pmatrix}
			K & K\\
			0 & K
		\end{pmatrix},
		\quad
		a1+bg\mapsto\begin{pmatrix}
			a+b & b\\
			0 & a+b
		\end{pmatrix},
	\]
	is an algebra isomorphism onto its image. 
	\end{enumerate}
\end{exercise}

If $A$ is an algebra over $K$ and $\rho\colon G\to \mathcal{U}(A)$
is a group homomorphism, where $\mathcal{U}(A)$ is the group of units of $A$, then 
the map \[
	K[G]\to A,\quad 
\sum_{g\in G}\lambda_gg\mapsto\sum_{g\in G}\lambda_g\rho(g),
\]
is an algebra homomorphism. 

\begin{exercise}
	Let $G=C_3$ be the (multiplicative) cyclic group of order three. Prove that
	$\R[G]\simeq\R\times\C$.
% 	Escribamos $G=\langle g:g^3=1\rangle$ y sea 
% 	\[
% 		\varphi\colon\R[G]\to\R\times\C,
% 		\quad
% 		g\mapsto (1,\omega),
% 	\]
% 	donde $\omega$ es una raíz cúbica primitiva de la unidad. Entonces
% 	$\varphi$ es inyectivo pues
% 	$0=\varphi(a1+bg+cg^2)=(a+b+c,a+b\omega+c\omega^2)$ implica que $a=b=c=0$.
% 	Luego $\varphi$ es un isomorfismo pues
% 	$\dim_\R\R[G]=\dim_\R(\R\times\C)=3$. 
\end{exercise}

\begin{exercise}
\label{xca:isos_dihedral}
	Let $G=\langle r,s:r^3=s^2=1,\,srs=r^{-1}\rangle$ be the dihedral group of six elements. 
	Prove the following statements:
	\begin{enumerate}
	    \item $\C[G]\simeq\C\times\C\times M_2(\C)$.
	    \item $\Q[G]\simeq\Q\times\Q\times M_2(\Q)$.
	\end{enumerate}  
% 	Sea $\omega$ una raíz cúbica de la unidad y sean  
% 	\[
% 		R=\begin{pmatrix}
% 			\omega & 0\\
% 			0 & \omega^2
% 		\end{pmatrix},
% 		\quad
% 		S=\begin{pmatrix}
% 			0 & 1\\
% 			1 & 0
% 		\end{pmatrix}.
% 	\]
% 	Un cálculo sencillo muestra que $R^2=S^2=I$ y que $SRS=R^{-1}$. Sea
% 	\[
% 		\varphi\colon\C[G]\to\C\times\C\times M_2(\C),\quad
% 		r\mapsto (1,1,R),\quad
% 		s\mapsto (1,-1,S).
% 	\]
% 	Es fácil ver que $\varphi$ es un morfismo de álgebras. Veamos que es
% 	biyectivo. Como $\dim_{\C}\C[G]=\dim_{\C}(\C\times\C\times M_2(\C))=6$,
% 	basta ver que $\varphi$ es inyectivo. Si 
% 	\[
% 		\alpha=a_0+a_1r+a_2r^2+(b_0+b_1r+b_2r^2)s\in\ker\varphi,
% 	\]
% 	entonces 
% 	\[
% 		0=\varphi(\alpha)=\left(u,v,\begin{pmatrix} \alpha_{11} & \alpha_{12}\\\alpha_{21}&\alpha_{22}\end{pmatrix}\right), 
% 	\]
% 	donde
% 	\begin{align*}
% 		&u = a_0+a_1+a_2+b_0+b_1+b_2, && v = a_0+a_1+a_2-b_0-b_1-b_2,\\
% 		&\alpha_{11}=a_0+a_1\omega+a_2\omega^2, && \alpha_{12}=b_0+b_1\omega+b_2\omega^2,\\
% 		&\alpha_{21}=b_0+b_2\omega+b_1\omega^2, && \alpha_{22}=a_0+a_2\omega+a_1\omega^2.
% 	\end{align*}
% 	Un cálculo sencillo muestra que estas ecuaciones implican que
% 	$\alpha=0$ y luego $\varphi$ es inyectiva.  
\end{exercise}



Maschke's theorem states that, if $G$ is a finite group, 
then the group algebra $\C[G]$ is semisimple. By Molien's theorem, 
\[
\C[G]\simeq \prod_{i=1}^k M_{n_i}(\C),
\]
where $k$ is the number of (isomorphism classes of) 
simple $\C[G]$-modules. Moreover, 
\[
|G|=\dim\C[G]=\sum_{i=1}^k n_i^2.
\]

\begin{theorem}
    Let $G$ be a finite group. The number of simple modules of $\C[G]$
    coincides with the number of conjugacy classes of $G$. 
\end{theorem}

\begin{proof}
    By Molien's theorem, $\C[G]\simeq\prod_{i=1}^kM_{n_i}(\C)$. Thus 
    \[
		Z(\C[G])\simeq\prod_{i=1}^kZ(M_{n_i}(\C))\simeq\C^k.
	\]
	In particular, $\dim Z(\C[G])=k$. If $\alpha=\sum_{g\in
	G}\lambda_gg\in Z(\C[G])$, then $h^{-1}\alpha h=\alpha$ for all $h\in
	G$. Thus 
	\[
		\sum_{g\in G}\lambda_{hgh^{-1}}g=
		\sum_{g\in G}\lambda_g h^{-1}gh=\sum_{g\in G}\lambda_gg
	\]
	and hence $\lambda_{g}=\lambda_{hgh^{-1}}$ for all $g,h\in G$. 
  Note that 
  coefficients constant on conjugacy classes give a central element. Thus 
  a basis for 
	$Z(\C[G])$ is given by elements of the form 
	\[
		\sum_{g\in K}g,
	\]
  where $K$ is a conjugacy class of $G$. Therefore $\dim Z(\C[G])$ is equal to
  the number of conjugacy classes of $G$.
\end{proof}

%\begin{project}{Which algebras are group algebras?}
\begin{optional}

\section{Which algebras are group algebras?}

\begin{example}
\label{exa:C4}
    Let $G=C_4$ be the cyclic group of order four. Then
    $G$ has four simple modules and 
    $\C[G]\simeq\C^4$. 
\end{example}

\begin{example}
\label{exa:S3}
    Let $G=\Sym_3$. Then $G$ has three simple modules and
    \[
    \C[G]\simeq\C\times\C\times M_2(\C).
    \]
\end{example}

\begin{open}[Brauer]
\index{Brauer's problem}
    Which complex algebras are group algebras of finite groups? 
\end{open}

The problem remains only partially understood. 
Examples \ref{exa:C4} and \ref{exa:S3} show
that $\C^4$ and $\C^2\times M_2(\C)$ are complex group algebras. 

\begin{xca}
    Is $\C^2\times M_2(\C)\times M_3(\C)$ a complex group algebra?  
\end{xca}

% No. Let $G$ be a group of order 15. Since groups of order 15 are abelian,
% $G$ has 15 conjugacy classes. 
\end{optional} 

\begin{graybox}
\project 
\section{The isomorphism problem for group algebras}
\label{project:isomorphism_problem}

\subsection{Group rings}

Let $R$ be a unitary commutative ring and $G$ be a group. Recall that
\emph{group ring} is a free $R$-module which is also a ring. As a free
$R$-module, $G$ is a basis for $R[G]$. Thus every element $\alpha\in R[G]$ can
be written uniquely as a sum of the form 
\[
\alpha=\sum_{g\in G}r_gg,\quad r_g\in R.
\]
Note that the elements of $R$ (more precisely, the elements of the form $r1_G$
for $r\in R$) are central in $R[G]$. 

\begin{xca}
\label{xca:augmentation_ideal}
    Let $R$ be a commutative ring and $G$ be a group. Let 
    \[
    f\colon R[G] \to R,\quad 
    \sum_{g \in G} r_g g\mapsto \sum_{g \in G} r_g.
    \]
    Prove the following statements:
    \begin{enumerate}
        \item $f$ is a surjective ring
            homomorphism. 
        \item $\{g - 1_G: 1_G\ne g \in G\}$ is a    basis for the free $R$-module $\ker f$.
    \end{enumerate}
\end{xca}

The kernel of the homomorphism of Exercise \ref{xca:augmentation_ideal} is
called the \emph{augmentation ideal} of the group ring. 
%We used this ideal
%before (e.g. Propositions \ref{pro:KGsemisimple} and Appendix
%\ref{section:local}). 

Note that $R[G]$ is a left $R$-module with
\[
\mu\left(\sum_{g\in G}\lambda_gg\right)=\sum_{g\in G}(\mu\lambda_g)g.
\]

\subsection{Hurewicz' theorem}

\begin{theorem}[Hurewicz]
    \label{thm:Hurewicz}
    \index{Hurewicz' theorem}
    Let $G$ be a group and $I$ be the augmentation ideal of $\Z[G]$. 
    Then $G/[G,G]\simeq I/I^2$ as (abelian) groups. 
\end{theorem}

\begin{proof}
    Let $\varphi\colon G\to I/I^2$, $g\mapsto g-1_G+I^2$. Since $g-1_G\in I$ for all $g\in G$, $\varphi$ is well-defined. The map $\varphi$ is a group homomorphism. Since 
    $(g-1_G)(h-1_G)\in I^2$, 
    \begin{align*}
    \varphi(gh) &= gh-1_G+I^2\\
    &=gh-1_G-(g-1_G)(h-1_G)+I^2+I^2\\
    &=g-1_G+h-1_G+I^2\\
    &=\varphi(g)+\varphi(h)
    \end{align*}
    holds for all $g,h\in G$. 

    Since $[G,G]\subseteq\ker\varphi$, there exists a group homomorphism
    \[
    \overline{\varphi}\colon G/[G,G]\to I/I^2,\quad 
    g[G,G]\mapsto g-1_G+I^2.
    \]
    We claim that $\overline{\varphi}$ is an isomorphism. 
    Let us construct the inverse of $\overline{\varphi}$. Let 
    \[
    \psi\colon I\to G/[G,G],\quad 
    \sum_{g\in G}m_g(g-1_G)\mapsto \left(\prod_{g\in G}g^{m_g}\right)[G,G].
    \]
    Since $G/[G,G]$ is abelian, the map $\psi$ is well-defined, that is
    the order of the factors in $\prod_{g\in G}g^{m_g}$ does not matter. Note that 
    $I^2\subseteq\ker\psi$, as 
    $\{(g-1_G)(h-1_G):g,h\in G\}$ generates the additive group $I^2$ 
    and 
    \begin{align*}
        \psi((g-1_G)(h-1_G))&=\psi( (gh-1_G)-(g-1_G)-(h-1_G))\\
        &=(ghg^{-1}h^{-1})[G,G]\\
        &=[G,G].
    \end{align*}
    Therefore there exists a group homomorphism
    \[
    \overline{\psi}\colon I/I^2\to G/[G,G],\quad 
    \sum_{g\in G}m_g(g-1_G)+I^2\mapsto \left(\prod_{g\in G}g^{m_g}\right)[G,G].
    \]
    A direct calculation shows that $\overline{\psi}$ is the inverse 
    of $\overline{\varphi}$. 
\end{proof}

\subsection{The isomorphism problem}\
%Recall that if $R$ is a unitary commutative ring 
%and $G$ is a group, then one defines the group ring $R[G]$. 
% More precisely,
% $R[G]$ is the set of finite linear combinations
% \[
%     \sum_{g\in G}\lambda_gg
% \]
% where $\lambda_g\in R$ and $\lambda_g=0$ for all but finitely many $g\in G$.
% One easily proves that $R[G]$ is a ring with
% addition
% \[
% \left(\sum_{g\in G}\lambda_gg\right)+\left(\sum_{g\in G}\mu_gg\right)
% =\sum_{g\in G}(\lambda_g+\mu_g)(g)
% \]
% and multiplication
% \[
% \left(\sum_{g\in G}\lambda_gg\right)\left(\sum_{h\in G}\mu_hh\right)
% =\sum_{g,h\in G}\lambda_g\mu_h(gh).
% \]
%Note that $R[G]$ is a left $R$-module with
%\[
%\mu\left(\sum_{g\in G}\lambda_gg\right)=\sum_{g\in G}(\mu\lambda_g)g.
%\]
%
%In this section, we will briefly discuss the
%following natural problem: 

\begin{que}[The isomorphism problem]
\label{question:iso}
\index{The isomorphism problem}
 Let $R$ be a unitary commutative ring and let 
 $G$ and $H$ be groups. Assume that $R[G]\simeq R[H]$ as
 unital $R$-algebras. Does $G\simeq H$?
\end{que}

For general information on Question \ref{question:iso} we refer to the survey
paper \cite{MR4472590}.

\begin{xca}
    Prove that if $G$ and $H$ are isomorphic groups, then $K[G]\simeq K[H]$.
\end{xca}

\begin{xca}
    Let $G$ and $H$ be groups. Prove that if $\Z[G]\simeq\Z[H]$, then
    $R[G]\simeq R[H]$ for any commutative ring $R$.
\end{xca}

The previous exercise suggests the importance of the following 
instance of Question \ref{question:iso}:

\begin{que}
\label{question:IP}
    Let $G$ and $H$ be groups. Does $\Z[G]\simeq\Z[H]$ imply $G\simeq H$?
\end{que}

For finite groups, the isomorphism problem has an affirmative answer in several
important cases, including abelian, metabelian, and nilpotent groups. However,
it is false in general. In 2001 Hertweck found a counterexample of order
$2^{21}97^{28}$, see \cite{MR1847590}.

\begin{que}[The modular isomorphism problem]
\label{question:MIP}
    Let $p$ be a prime number. Let
    $G$ and $H$ be finite $p$-groups and let $K$ be a field of characteristic $p$.
    Does $K[G]\simeq K[H]$ imply $G\simeq H$?
\end{que}

Question \ref{question:MIP} has an affirmative answer in several cases.
However, this is not true in general. This question was answered by
Garc\'ia-Lucas, Margolis and del R\'io \cite{MR4373245}. They found two
non-isomorphic groups $G$ and $H$ both of order $512$ such that $K[G]\simeq
K[H]$ for every field $K$ of characteristic two.
\end{graybox}

\section{Rickart's theorem}

We now consider Jacobson's semisimplicity problem. 

\begin{question}
\label{Jacobson's semisimplicity problem}
Let $G$ be a group and $K$ be a field. When $J(K[G])=\{0\}$?
\end{question}

As an application of Amitsur's theorem \ref{thm:Amitsur}, 
we prove that 
complex group algebras have null Jacobson radical.
This is known as 
Rickart's theorem. The original proof found by Rickart 
uses complex analysis. Here, however, 
we present an algebraic proof. 

\begin{theorem}[Rickart]
\index{Rickart's theorem}
\label{thm:Rickart}
    Let $G$ be a group. Then $J(\C[G])=\{0\}$.
\end{theorem}

To prove the theorem, we need a lemma.

\begin{lemma}
Let $G$ be a group. Then $J(\C[G])$ is nil.        
\end{lemma}

\begin{proof}
    We need to show that every element of $J(\C[G])$ is nilpotent. 
    If $G$ is countable, then the result follows from Amitsur's theorem \ref{thm:Amitsur}. So assume that 
    $G$ is not countable. Let $\alpha\in J(\C[G])$, say
    \[
    \alpha=\sum_{i=1}^n\lambda_ig_i,
    \]
    where $\lambda_1,\dots,\lambda_n\in\C$ and $g_1,\dots,g_n\in G$. Let $H=\langle g_1,\dots,g_n\rangle$.
    Then $\alpha\in \C[H]$ and $H$ is countable. We claim that $\alpha\in J(\C[H])$. Decompose
    $G$ as a disjoint union 
    \[
    G=\bigcup_\lambda x_\lambda H
    \]
    of cosets of $H$ in $G$. Then $\C[G]=\bigoplus_\lambda x_\lambda\C[H]$ and
    hence $\C[G]=\C[H]\oplus K$ for some right module $K$ over $\C[H]$ (this follows
    from the fact that one of the cosets is that of $H$). Since $\alpha\in J(\C[G])$, for each 
    $\beta\in\C[H]$ there exists $\gamma\in\C[G]$ such that 
    $\gamma(1-\beta\alpha)=1$. Write $\gamma=\gamma_1+\kappa$ for $\gamma_1\in\C[H]$ and $\kappa\in K$. Then
    \[
    1=\gamma(1-\beta\alpha)=\gamma_1(1-\beta\alpha)+\kappa(1-\beta\alpha)
    \]
    and hence $\kappa(1-\beta\alpha)\in K\cap \C[H]=\{0\}$, as $\beta\in\mathbb{C}[H]$. 
    Since $1=\gamma_1(1-\beta\alpha)$, it follows that
    $\alpha\in J(\C[H])$ and the lemma follows from Amitsur's theorem \ref{thm:Amitsur}.  
\end{proof}

We now prove the theorem. 

\begin{proof}[Proof of Theorem \ref{thm:Rickart}]
    For $\alpha=\sum_{i=1}^n\lambda_ig_i\in\C[G]$ let 
    \[
    \alpha^*=\sum_{i=1}^n\overline{\lambda_i}g_i^{-1}.
    \]
    Then $\alpha\alpha^*=0$ if and only if $\alpha=0$ and, moreover, 
    $(\alpha\beta)^*=\beta^*\alpha^*$ for all $\beta\in\C[G]$. 
    Assume that $J(\C[G])\ne\{0\}$ and let $\alpha\in J(\C[G])\setminus\{0\}$. Then
    $\beta=\alpha\alpha^*\in J(\C[G])$, as $J(\C[G])$ is an ideal of $\C[G]$. Moreover, the previous 
    lemma implies that $\beta$ is nilpotent. Note that $\beta\ne 0$, as $\alpha\ne0$. Since $\beta^*=\beta$,  
    \[
    (\beta^m)^*=(\beta^*)^m=\beta^m
    \]
    for all $m\geq1$. If there exists $k\geq2$ such that $\beta^k=0$ and $\beta^{k-1}\ne 0$, then
    \[
    \beta^{k-1}\left(\beta^{k-1}\right)^*=\beta^{2k-2}=0
    \]
    and hence $\beta^{k-1}=0$, a contradiction. Thus $\beta=0$ and therefore $\alpha=0$. 
\end{proof}

\begin{exercise}
	If $G$ is a group, then $J(\R[G])=0$. 
\end{exercise}

% To obtain a consequence of Rickart's theorem we need two lemmas. 

% \begin{lemma}[Nakayama]
% 	\label{lem:Nakayama}
% 	\index{Nakayama's lemma}
% 	Let $R$ be a unitary ring and $M$ be a finitely generated module. If 
% 	$J(R)\cdot M=M$, then $M=\{0\}$.
% \end{lemma}

% \begin{proof}
%     Since $M$ is finitely generated, we may assume that 
% 	$M=(x_1,\dots,x_n)$. Since $x_n\in M=J(R)\cdot M$, 
% 	there exist $r_1,\dots,r_n\in J(R)$ such that $x_n=r_1\cdot x_1+\cdots+r_n\cdot x_n$, that is 
% 	$(1-r_n)\cdot x_n=\sum_{j=1}^{n-1}r_j\cdot x_j$. 
% 	Since $1-r_n$ is invertible, there exists $s\in R$ such that $s(1-r_n)=1$. Thus 
% 	$x_n=\sum_{j=1}^{n-1}(sr_j)\cdot x_j$ 
% 	and hence $M=(x_1,\dots,x_{n-1})$. Repeating this procedure several times 
% 	one obtains $M=\{0\}$.
% \end{proof}

% \begin{lemma}
% 	\label{lem:Rickart}
% 	Let $\iota\colon R\to S$ be a homomorphism of unitary rings. If	
% 	\[
% 	S=\iota(R)x_1+\cdots+\iota(R)x_n,
% 	\]
% 	where each $x_j$ is such that $x_jy=yx_j$ for all $y\in\iota(R)$, then 
% 	$\iota(J(R))\subseteq J(S)$.
% \end{lemma}

% \begin{proof}
% 	We claim that $J=\iota(J(R))$ acts trivially on each simple $S$-module $M$.
% 	If is $M$ is a simple module over $S$, then, in particular, $M=S\cdot m$ for some $m\ne0$. 
% 	Now $M$ is a module over $R$ with $r\cdot m=\iota(r)\cdot m$. Since 
% 	\[
% 		M=S\cdot m=(\iota(R)x_1+\cdots+\iota(R)x_n)\cdot m=\iota(R)\cdot (x_1\cdot m)+\cdots+\iota(R)\cdot (x_n\cdot m),
% 	\]
% 	it follows that 
% 	$M$ is finitely generated as a module over $\iota(R)$. Moreover, 
% 	\[
% 	J(R)\cdot
% 	M=J\cdot M=\iota(J)\cdot M
% 	\]
% 	is an $S$-submodule of $M$, as 
% 	\[
% 		x_j\cdot (J\cdot M)=(x_j J)\cdot M=(J x_j)\cdot M=J\cdot (x_j\cdot M)\subseteq J\cdot M.
% 	\]
% 	Since $M\ne\{0\}$, Nakayama's lemma implies that $J(R)\cdot M\subsetneq M$. The simplicity of 
% 	the $S$-module $M$ implies that $J(R)\cdot M=\{0\}$.
% \end{proof}

% We now obtain the following consequence of Rickart's theorem. 

% \begin{theorem}
% 	If $G$ is a group, then $J(\R[G])=0$. 
% \end{theorem}

% \begin{proof}
% 	Let $\iota\colon \R[G]\to\C[G]$ be the canonical inclusion. Since 
% 	\[
% 	\C[G]=\R[G]+i\R[G],
% 	\]
% 	Lemma~\ref{lem:Rickart} and Rickart's theorem imply that 
% 	$\iota(J(\R[G]))\subseteq J(\C[G])=0$. Thus $J(\R[G])=0$, as $\iota$ is injective. 
% \end{proof}


\begin{definition}
	\index{Ring!semiprime}
	A ring $R$ \emph{semiprime} if 
	$aRa=\{0\}$ implies $a=0$.
\end{definition}

\begin{proposition}
	Let $R$ be a ring. The following statements are equivalent: 
	\begin{enumerate}
		\item $R$ is semiprime.
		\item If $I$ is a left ideal such that $I^2=\{0\}$, then $I=\{0\}$.
		\item If $I$ is an ideal such that $I^2=\{0\}$, then $I=\{0\}$.
		\item $R$ does not contain non-zero nilpotent ideals.
	\end{enumerate}
\end{proposition}

\begin{proof}
	We first prove that $1)\implies2)$. If $I^2=\{0\}$ y $x\in I$, then
	$xRx\subseteq I^2=\{0\}$ and thus $x=0$. 
    
    The implications $2)\implies3)$
	and $4)\implies3)$ are both trivial. 
    
    Let us prove that $3)\implies4)$.  If
	$I$ is a non-zero nilpotent ideal, let $n\in\Z_{>0}$ be minimal such that
	$I^n=\{0\}$.  Since $(I^{n-1})^2=\{0\}$, it follows that $I^{n-1}=\{0\}$, a
	contradiction.  
    
    Finally, we prove that $3)\implies1)$. Let $a\in R$ be such
	that $aRa=\{0\}$. Then $I=RaR$ is an ideal of $R$ such that $I^2=\{0\}$. Thus 
	$RaR=\{0\}$. This means that $Ra$ and $aR$ are ideals such that
	$(Ra)R=R(aR)=\{0\}$ (for example, $R(aR)\subseteq RaR=\{0\}\subseteq aR$). 
	Moreover, since $(Ra)(Ra)=\{0\}$ and $(aR)(aR)=\{0\}$, it follows that
	$aR=Ra=\{0\}$. 
	This implies that $\Z a$ is an ideal of $R$, as $R(\Z a)\subseteq \Z(Ra)=\{0\}$ and 
	$(\Z a)R\subseteq aR=\{0\}$. Now $(\Z a)(\Z a)\subseteq (\Z a)R=\{0\}$ and hence
	$a=0$, as $\Z a=\{0\}$. 
\end{proof}

Two consequences:

\begin{exercise}
	A commutative ring is semiprime if and only if it does not contain non-zero
	nilpotent elements. 
\end{exercise}

% tomar $\alpha$ tal que $\alpha^2=0$ y sea $A=K[G]\alpha$. Como $A^2=0$, $A=0$ y entonces $\alpha=0$.

\begin{exercise}
\label{xca:D_semiprime_semiprimitive}
	Let $D$ be a division ring. 
	\begin{enumerate}
		\item $D[X]$ is semiprime and semiprimitive.
		\item $D[\![X]\!]$ is semiprime and it is not semiprimitive.
	\end{enumerate}
\end{exercise}

% We will prove 
% in Lecture \ref{09} (Corollary \ref{cor:C[G]_semiprime}) 
% that the if $G$ is a group, 
% then the ring $\C[G]$ is semiprime. 

\begin{corollary}
\label{cor:C[G]_semiprime}
	The ring $\C[G]$ is semiprime.
\end{corollary}

\begin{proof}
	Since $J(\C[G])=\{0\}$ by Rickart's theorem and the Jacobson radical
	contains every nil ideal by Proposition~\ref{pro:nilJ}, it follows that
	$\C[G]$ does not contain non-trivial nil ideals. Thus $\C[G]$ does not
	contain non-trivial nilpotent ideals and hence $\C[G]$ is semiprime.
\end{proof}

\begin{exercise}
	Prove that $Z(\C[G])$ is semiprime.
\end{exercise}

We now characterize when complex group algebras 
are left artinian. For that purpose,
we need a lemma. This is similar to one of the implications proved in Proposition \ref{pro:semisimple}. However,
in the arbitrary setting we are considering, we need to use Zorn's lemma. 

\begin{lemma}
\label{lem:direct_summand}
    Let $M$ be a semisimple module and $N$ be a submodule. 
    Then $N$ is a direct summand.
\end{lemma}

\begin{proof}[Sketch of the proof]
    Let $M=\oplus_{i\in I}M_i$ be a direct sum of simple submodules. 
    Since each $N\cap M_i$ is a submodule of $M_i$ and $M_i$ is simple, it follows
    that $N\cap M_i=\{0\}$ or $N\cap M_i=M_i$. If
    $N\cap M_i=M_i$ for all $i\in I$, then $N=M$ and the lemma is proved. So we may assume
    that there exists $i\in I$ such that $N\cap M_i=\{0\}$. Let $X$ be the set
    of subsets $J$ of $I$ such that $N\cap (\oplus_{j\in J}M_j)=\{0\}$. Our assumptions
    imply that $X$ is non-empty. Zorn's lemma implies the existence of 
    a maximal element $K$. Let $N_1=\bigoplus_{k\in K}M_k$. We claim that
    $N\bigoplus N_1=M$. If not, there exists $i\in I$ such that
    $M_i\not\subseteq N\oplus N_1$. The simplicity of $M_i$ implies that
    $M_i\cap (N\oplus N_1)=\{0\}$, which contradicts the maximality of $K$. 
\end{proof}

A direct application of the lemma proves that
complex group algebras of infinite groups are never semisimple. 

\begin{proposition}
    \label{pro:KGsemisimple}
    If $G$ is an infinite group, then $\C[G]$ is not semisimple. 
\end{proposition}

\begin{proof}
	Assume that $R=\C[G]$ is semisimple.  Let $I$ 
	be the augmentation ideal of $R$, that is
	\[
	I=\left\{\alpha=\sum_{g\in G}\lambda_gg\in R:\sum_{g\in G}\lambda_g=0\right\}.
	\]
	By Lemma~\ref{lem:direct_summand}, 
	there exists a non-zero left ideal $J$ such that 
	$R=I\oplus J$. Since $R$ is unitary, there exist $e\in I$ and $f\in J$ such that
	$1=e+f$. If
	$x\in I$, then $x=xe+xf$ and hence $xf=x-xe\in I\cap J=\{0\}$. Since 
	$x=xe$ for all $x\in I$, it follows that $e=e^2$. Similarly, one proves
	that $f^2=f$. Moreover, $ef=0$, as $ef\in I\cap J=\{0\}$.  Since $I$ 
	is the augmentation ideal of $R$ and $If=(Re)f=R(ef)=\{0\}$ 
 (note that $I=Re$ because $x=xe$ for all $x\in I$), we conclude that
	$(g-1_G)f=0$
	for all $g\in G$, as $g-1_G\in I$ for all $g\in G$. If $f=\sum_{h\in
	G}\lambda_hh$ (finite sum) and $g\in G$, then  
	\[
	f=gf=\sum_{h\in G}\lambda_h(gh)=\sum_{h\in
	G}\lambda_{g^{-1}h}h.
	\]
        Thus $\lambda_h=\lambda_{g^{-1}h}$ for all $g,h\in G$. Since $G$ 
	is infinite, some $\lambda_g=0$ and hence $f=0$. Thus $e=1$ and $I=\C[G]$, a contradiction. 
% 	If $f=0$, then $e=1$ and $I=\C[G]$, a contradiction.  
% 	, a contradiction because 
% 	$f\ne 0$ implies that the sum that defines $f$ should be an infinite sum.
\end{proof}

% The ideal $I(G)$ used in the proof of the previous proposition 
% is known as the \emph{augmentation ideal} 
% of $\C[G]$.

\begin{theorem}
	Let $G$ be a group. Then $\C[G]$ 
	is left artinian if and only if 
	$G$ is finite. 
\end{theorem}

\begin{proof}
    If $G$ is finite, then $\C[G]$ is left artinian because $\dim\C[G]=|G|<\infty$. So assume that 
    $G$ is infinite. By Rickart's theorem,   
	$J(\C[G])=0$. Moreover, $\C[G]$
	is not semisimple by the previous proposition. Thus
	$\C[G]$ is not left artinian by Theorem~\ref{thm:SSartin=J}.
\end{proof}

\begin{graybox}
  \project 
  \section{Another proof of Rickart's theorem}

  We will present a different proof 
  of Rickart's theorem \ref{thm:Rickart}. 

  \begin{definition}
    \index{Ring!with an involution}
    \index{Involution}
    Let $R$ be a ring. An \emph{involution} 
    in $R$ is an additive map  
    $R\to R$, $x\mapsto x^*$, such that $x^{**}=x$ and  $(xy)^*=y^*x^*$ for all 
    $x,y\in R$.
  \end{definition}

  It follows that if $R$ is a unitary ring, then 
  $1^*=1$.

  \begin{example}
    The conjugation $z\mapsto\overline{z}$ is an 
    involution of $\C$.
  \end{example}

  \begin{example}
    The transpose map $X\mapsto X^T$ is an involution of $M_n(K)$.
  \end{example}

  \begin{example}
    Let $G$ be a group. 
    Then \[
      \left(\sum_{g\in G}\alpha_gg\right)^*=\sum_{g\in G}\overline{\alpha_g}g^{-1}
    \]
    is an involution of the group ring $\C[G]$.
  \end{example}

  For a group $G$, one defines the \emph{trace} 
  of an element 
  \[
    \alpha=\sum_{g\in G}\alpha_gg\in K[G]
  \]
  as $\trace(\alpha)=\alpha_1$. One proves that 
  the map 
  \[
    \trace\colon K[G]\to K,\quad 
    \alpha\mapsto\trace(\alpha)
  \]
  is $K$-linear such that   
  $\trace(\alpha\beta)=\trace(\beta\alpha)$ for 
  all $\alpha,\beta\in K[G]$. 

  \begin{xca}
    Let $G$ be a finite group and $K$ 
    be a field of characteristic not dividing the 
    order of $G$.
    Prove the following statements: 
    \begin{enumerate}
      \item If $\alpha\in K[G]$ is nilpotent, then $\trace(\alpha)=0$.
      \item If $\alpha\in K[G]$ is idempotent, then $\trace(\alpha)=\dim
        (K[G]\alpha)/|G|$.
    \end{enumerate}
  \end{xca}

  \begin{xca}
    Prove that  
    $\langle\alpha,\beta\rangle=\trace(\alpha\beta^*)$, $\alpha,\beta\in\C[G]$, 
    defines an inner product in $\C[G]$.
  \end{xca}

  \begin{lemma}
    \label{lem:algebraico}
    Let $G$ be a group. Prove that if $J(\C[G])\ne 0$, then there exists $\alpha\in J(\C[G])$ such that  
    $\trace(\alpha^{2^m})\in\R_{\geq1}$ 
    for all $m\geq1$.
  \end{lemma}

  \begin{proof}
    Let $\alpha=\sum_{g\in G}\alpha_gg\in\C[G]$. Then 
    \[
      \trace(\alpha^*\alpha)
      =\sum_{g\in G}\overline{\alpha_g}\alpha_g
      =\sum_{g\in G}|\alpha_g|^2\geq|\alpha_1|^2
      =|\trace(\alpha)|^2.
    \]
    By induction, one proves that, if 
    $\alpha$ is such that $\alpha^*=\alpha$, then  
    \[
      \trace(\alpha^{2^m})\geq|\trace(\alpha)|^{2^m}
    \]
    for all $m\geq1$. 

    Let $\beta=\sum_{g\in G}\beta_gg\in J(\C[G])$ be such that $\beta\ne0$. Since 
    $\trace(\beta^*\beta)=\sum_{g\in G}|\beta_g|^2\ne0$ and $J(\C[G])$ is an ideal, 
    \[
      \alpha=\frac{\beta^*\beta}{\trace(\beta^*\beta)}\in J(\C[G]).
    \]
    Then $\alpha$ is such that $\alpha^*=\alpha$ and $\trace(\alpha)=1$.
    Hence $\trace(\alpha^{2^m})\geq 1$  for all $m\geq1$.
  \end{proof}

  Exercise~\ref{exa:norma} implies that $\C[G]$ with 
  $\operatorname{dist}(\alpha,\beta)=|\alpha-\beta|$ is a metric space. In this metric space, the map 
  $\C[G]\to\C$, $\alpha\mapsto \trace(\alpha)$, is a continuous map. 

  \begin{lemma}
    \label{lem:phi_diferenciable}
    Let $\alpha\in J(\C[G])$. The map 
    \[
      \varphi\colon\C\to\C[G],\quad
      \varphi(z)=(1-z\alpha)^{-1},
    \]
    is continuous, differentiable and $\varphi(z)=\sum_{n\geq0}\alpha^nz^n\in\C[G]$ if  $|z|$
    is sufficiently small. 
  \end{lemma}

  \begin{proof}	
    Let $y,z\in\C$. Since $\varphi(y)$ and $\varphi(z)$ commute, 
    \begin{equation}
      \label{eq:Rickart}
      \begin{aligned}
        \varphi(y)-\varphi(z)&=\left( (1-z\alpha)-(1-y\alpha)\right)(1-y\alpha)^{-1}(1-z\alpha)^{-1}\\
                             &=(y-z)\alpha\varphi(y)\varphi(z).
      \end{aligned}
    \end{equation}
    Hence $|\varphi(y)|\leq|\varphi(z)|+|y-z||\alpha\varphi(y)||\varphi(z)|$ and therefore 
    \[
      |\varphi(y)|\left( 1-|y-z||\alpha\varphi(z)|\right)\leq|\varphi(z)|.
    \]
    Fix $z$ and choose $y$ sufficiently close to $z$ in such a way that $1-|y-z||\alpha\varphi(z)|\geq1/2$. Then
    $|\varphi(y)|\leq2|\varphi(z)|$. Using~\eqref{eq:Rickart} one obtains that 
    $|\varphi(y)-\varphi(z)|\leq2|y-z||\alpha||\varphi(z)|^2$. 
    Hence $\varphi$ is continuous. By~\eqref{eq:Rickart}, 
    \[
      \varphi'(z)
      =\lim_{y\to z}\frac{\varphi(y)-\varphi(z)}{y-z}
      =\lim_{y\to z}\alpha\varphi(y)\varphi(z)
      =\alpha\varphi(z)^2
    \]
    for all $z\in\C$.

    If $z$ is such that $|z||\alpha|=|z\alpha|<1$, then  
    \[
      \varphi(z)-\sum_{n=0}^Nz^n\alpha^n
      =\varphi(z)\left(1-(1-z\alpha)\sum_{n=0}^Nz^n\alpha^n\right)
      =\varphi(z)(z\alpha)^{N+1}.
    \]
    Then 
    \[
      \left|\varphi(z)-\sum_{n=0}^Nz^n\alpha^n\right|\leq|\varphi(z)||z\alpha|^{N+1}.
    \]
    Since $\varphi(z)$ is bounded close to $z=0$, we conclude that 
    \[
      \left|\varphi(z)-\sum_{n=0}^Nz^n\alpha^n\right|\to0
    \]
    if $N\to\infty$.
  \end{proof}

  We now provide an alternative proof of Rickart's theorem: 

  \begin{theorem}[Rickart]
    \index{Rickart's theorem}
    If $G$ is a group, then $J(\C[G])=\{0\}$.
  \end{theorem}

  \begin{proof}
    Let $\alpha\in J(\C[G])$ and $\varphi(z)=(1-\alpha z)^{-1}$. Let 
    \[
      f\colon\C\to \C,\quad  
      f(z)=\trace\varphi(z)=\trace\left((1-z\alpha)^{-1}\right).
    \]
    By Lemma~\ref{lem:phi_diferenciable},
    $f(z)$ is an entire function such that  $f'(z)=\trace(\alpha\varphi(z)^2)$ and 
    \begin{equation}
      \label{eq:Taylor}
      f(z)=\sum_{n=0}^\infty z^n\trace(\alpha^n)
    \end{equation}
    if $|z|$ is sufficiently small. In particular,~\eqref{eq:Taylor} is the Taylor series of $f(z)$ around the origin. This implies that the series converges to $f(z)$ for all $z\in\C$. In particular,
    \begin{equation}
      \label{eq:limite}
      \lim_{n\to\infty}\trace(\alpha^n)=0.
    \end{equation}
    On the other hand, if $\alpha\ne0$, Lemma~\ref{lem:algebraico} implies that 
    $\trace(\alpha^{2^m})\geq1$ for all $m\geq0$. This contradicts the limit computed in~\eqref{eq:limite}. Hence $\alpha=0$.
  \end{proof}

  For a corollary, we need a consequence of 
  Nakayama's lemma \ref{lem:Nakayama}.


  \begin{lemma}
    \label{lem:Rickart}
    Let $R$ and $S$ be unitary rings and 
    $\iota\colon R\to S$ be a ring homomorphism. If  
    \[
      S=\iota(R)x_1+\cdots+\iota(R)x_n,
    \]
    where each $x_j$ is such that
    $x_jy=yx_j$ for all $y\in\iota(R)$, then
    $\iota(J(R))\subseteq J(S)$.
  \end{lemma}

  \begin{proof}
    We claim that $J=\iota(J(R))$ acts trivially on each simple $S$-module $M$.
    Let $M$ be a simple $S$-module. Write 
    $M=Sm$ for some $m\ne0$. Then $M$ is 
    an $R$-module with $r\cdot m=\iota(r)m$. Since 
    \[
      M=Sm=(\iota(R)x_1+\cdots+\iota(R)x_n)m=\iota(R)(x_1m)+\cdots+\iota(R)(x_nm),
    \]
    $M$ is finitely generated as an $\iota(R)$-module. Moreover, $J(R)\cdot
    M=JM=\iota(J)M$ is an $S$-submodule of $M$, as 
    \[
      x_j(JM)=(x_jJ)M=(Jx_j)M=J(x_jM)\subseteq JM.
    \]
    Since $M\ne\{0\}$, Nakayama's lemma implies that $J(R)\cdot M\subsetneq M$. Since 
    $M$ is a simple $S$-module, we conclude that 
    $J(R)M=\{0\}$.
  \end{proof}

  \begin{corollary}
    If $G$ is a group, then $J(\R[G])=\{0\}$. 
  \end{corollary}

  \begin{proof}
    Let $\iota\colon \R[G]\to\C[G]$ be the canonical inclusion. Since 
    \[
      \C[G]=\R[G]+i\R[G],
    \]
    Now Lemma \ref{lem:Rickart} and Rickart's theorem imply that 
    $\iota(J(\R[G]))\subseteq J(\C[G])=\{0\}$. We conclude that $J(\R[G])=0$, as the map $\iota$ is injective. 
  \end{proof}
\end{graybox}


\section{Maschke's theorem}

We now present another instance of the semisimplicity problem.
In this case, our result is for finite groups. 

\begin{theorem}[Maschke]
\index{Maschke's theorem}
	Let $G$ be a finite group. Then $J(K[G])=\{0\}$ if and only 
	if the characteristic of $K$ is zero 
	or does not divide the order of $G$. 
\end{theorem}

\begin{proof}
	Assume that $G=\{g_1,\dots,g_n\}$, where $g_1=1$. Let 
	\[
	\rho\colon K[G]\to K,
	\quad
	\alpha\mapsto\trace(L_{\alpha}),
	\]
	where 
	$L_{\alpha}(\beta)=\alpha\beta$. Then 
	\[
	\rho(g_i)=\begin{cases}
	    n & \text{if $i=1$,}\\
	    0 & \text{if $2\leq i\leq n$},
	\end{cases}
	\]
	as $L_{g_i}(g_j)=g_{i}g_j\ne g_j$ if $i\ne j$ and hence the matrix of 
	$L_{g_i}$ in the basis $\{g_1,\dots,g_n\}$ contains zeros in the main diagonal.

	Assume that $J=J(K[G])$ is non-zero and let 
	$\alpha=\sum_{i=1}^n\lambda_ig_i\in J\setminus\{0\}$. Without loss of generality
	we may assume that $\lambda_1\ne 0$ (if $\lambda_1=0$ there exists some 
	$\lambda_i\ne 0$ and we need to take $g_i^{-1}\alpha\in J$). Then 
	\[
		\rho(\alpha)=\sum_{i=1}^n \lambda_i\rho(g_i)=n\lambda_1.
	\]
	Since $G$ is finite, $K[G]$ is a finite-dimensional algebra and hence 
	$K[G]$ is left artinian. Since $J$ is a nilpotent ideal, 
	in particular, $\alpha$ is a nilpotent element. Then 
	$L_{\alpha}$ is nilpotent and hence $0=\rho(\alpha)=n\lambda_1$. This implies that
	the characteristic of the field $K$ divides $n$. 

	Conversely, let $K$ be a field of prime characteristic and that this prime divides 
	$n$. Let $\alpha=\sum_{i=1}^ng_i$. Since $\alpha
	g_j=g_j\alpha=\alpha$ for all $j\in\{1,\dots,n\}$, the set 
	$I=K[G]\alpha$ is an ideal of $K[G]$. Since, moreover,   
	\[
		\alpha^2=\sum_{i=1}^n g_i\alpha=n\alpha=0
	\]
	in the field $K$, it follows that $I$ is a nilpotent non-zero ideal. Thus $J(K[G])\ne\{0\}$, 
	as Proposition~\ref{pro:nilJ} yields $I\subseteq J(K[G])$.
\end{proof}

Since the Jacobson radical of a group algebra of a finite group contains 
every nil left ideal, the following consequence of the theorem follows immediately:

\begin{corollary}
	\label{cor:GfinitoNOnil}
	Let $G$ be a finite group. Then $K[G]$ does not contain non-zero nil left ideals. 
\end{corollary}


% \begin{proof}
% 	Es consecuencia inmediata del teorema de Maschke ya que $J(K[G])$ contiene a
% 	todo ideal a izquierda nil.	
% \end{proof}

%\index{Anillo!semisimple}
%Recordemos que un anillo unitario $R$ se dice \emph{semisimple} si para cada
%ideal $I$ de $R$ existe un ideal $J$ de $R$ tal que $R=I\oplus J$.
%
%%\begin{corollary}
%%	Sea $G$ un grupo finito y $K$ un cuerpo de característica coprima con el
%%	orden de $G$. Entonces $K[G]$ es semisimple.
%%\end{corollary}
%%
%%\begin{proof}
%%	
%%\end{proof}
%
%\begin{theorem}
%	Si $G$ es un grupo infinito, entonces $K[G]$ nunca es semisimple.
%\end{theorem}
%
%\begin{proof}
%	Sea $R=K[G]$ y supongamos que $R$ es semisimple.  Si $I$ es el ideal de
%	aumentación de $R$, existe un ideal no nulo $J$ de $R$ tal que $R=I\oplus
%	J$. Como $R$ es unitario, existen $e\in I$, $f\in J$ tales que $1=e+f$. Si
%	$x\in I$, entonces $x=xe+xf$ y luego $xf=x-xe\in I\cap J=\{0\}$. Como
%	entonces $x=xe$ para todo $x\in I$, en particular $e_1=e_1^2$. Análogamente
%	vemos que $e_2^2=e_2$. Además $ef=0$ pues $ef\in I\cap J=\{0\}$.  Como $I$
%	es el ideal de aumentación y $If=(Re)f=R(ef)=0$, se concluye que $(g-1)f=0$
%	para todo $g\in G$ pues $g-1\in I$. Si suponemos que $f=\sum_{h\in
%	G}\lambda_hh$, entonces 
%	\[
%	f=gf=\sum_{h\in G}\lambda_h(gh)=\sum_{h\in
%	G}\lambda_{g^{-1}h}h.
%	\]
%	Luego $\lambda_h=\lambda_{g^{-1}h}$ para todo $g,h\in G$, una contradicción
%	pues como $f\ne 0$ la suma que define a $f$ es infinita. 
%\end{proof}


\section{Herstein's theorem}

Our aim now is to answer the following question: When
a group algebra is algebraic? Herstein's theorem provides
a solution in the case of fields of characteristic zero. In prime characteristic,
the problem is still open. 

\begin{definition}
\index{Group!locally finite}
	A group $G$ is \emph{locally finite} if every finitely generated 
	subgroup of $G$ is finite. 
\end{definition}

If $G$ is a locally finite group, then every element $g\in G$ has finite order, as
the subgroup $\langle g\rangle$ is finite because it is finitely generated.

\begin{example}
    Every finite group is locally finite
\end{example}

\begin{example}
    The group $\Z$ is not locally finite because it is torsion-free.
\end{example}

\begin{example}
\index{Pr\"ufer's group}
	Let $p$ be a prime number. 
	The \emph{Pr\"ufer's group}  
	\[
		\Z(p^{\infty})=\{z\in\C:z^{p^n}=1\text{ for some $n\in\Z_{>0}$}\}, 
	\]
	is locally finite. 
\end{example}

\begin{example}
	Let $X$ be an infinite set and $\Sym_X$ be the set of bijective maps $X\to
	X$ moving only finitely many elements of $X$. Then 
	$\Sym_X$ is locally finite.
\end{example}

\index{Group!torsion}
A group $G$ is a \emph{torsion} group if every element of $G$
has finite order. Locally finite groups are torsion groups. 

\begin{example}
    Abelian torsion groups are locally finite. Let $G$ be a locally finite abelian group 
    and $H$ be a finitely generated subgroup. Since $G$ is an abelian torsion group, so is $H$. Thus
    $H$ is finite by the structure theorem of abelian groups. 
\end{example}

\begin{proposition}
\label{pro:exact_LI}
	Let $G$ be a group and $N$ be a normal subgroup of $G$. If $N$ and $G/N$
	are locally finite, then $G$ is locally finite.
\end{proposition}

\begin{proof}
	Let $\pi\colon G\to G/N$ be the canonical map and $\{g_1,\dots,g_n\}$ be a finite subset of $G$. 
	Since $G/N$ is locally finite, the subgroup $Q$ of $G/N$ generated by 
	$\pi(g_1),\dots,\pi(g_n)$ is finite, say
	\[
		Q=\{\pi(g_1),\dots,\pi(g_n),\pi(g_{n+1}),\dots,\pi(g_m)\}
	\]
	for some $g_{n+1},\dots,g_m\in G$. 
	
	For each $i,j\in\{1,\dots,n\}$ there exist $u_{ij}\in N$ and 
	$k\in\{1,\dots,m\}$ such that \[
	g_ig_j=u_{ij}g_k.
	\]
	Let $U$ be the subgroup of $N$
	generated by $\{u_{ij}:1\leq i,j\leq n\}$. Since $N$ is locally finite, $U$ is finite. Moreover, since 
	each $g_ig_jg_l$ can be written as 
	\[
		g_ig_jg_l=u_{ij}g_kg_l=u_{ij}u_{kl}g_t=ug_t
	\]
	for some $u\in U$ and $t\in\{1,\dots,m\}$, it follows that the subgroup 
	$H$ of $G$ generated by $\{g_1,\dots,g_n\}$ is finite, as 
	$|H|\leq m|U|$. 
\end{proof}

% \index{Group!solvable}
% A group $G$ is
% \emph{solvable} if there exists a sequence
% of subgroups 
% \begin{equation}
% 	\label{eq:resoluble}
% 	\{1\}=G_0\subsetneq G_1\subsetneq \cdots\subsetneq G_n=G
% \end{equation}
% where each $G_i$ is normal in $G_{i+1}$ and each 
% quotient $G_i/G_{i-1}$ is
% abelian.

% \begin{example}
%     Abelian groups are solvable. 
% \end{example}

% Subgroups and quotients of solvable groups are solvable. 

% \begin{example}
%     Groups of order $<60$ are solvable.
% \end{example}

% \begin{example}
%     $\Alt_5$ and $\Sym_5$ are not solvable. 
% \end{example}

% \index{Burnside's theorem}
% \index{Feit--Thompson's theorem}
% A famous theorem of Burnside states that 
% groups of order $p^aq^b$ for prime numbers $p$ and $q$ are solvable 
% A much harder theorem proved by Feit and Thompson states that
% groups of odd order are solvable.

% \begin{proposition}
% 	If $G$ is a solvable torsion group, 
% 	then $G$ is locally finite. 
% \end{proposition}

% \begin{proof}
% 	We proceed by induction on $n$, the length of the sequence~\eqref{eq:resoluble}. 
% 	If $n=1$, then $G$ is finite because it is abelian and a torsion group.
% 	Now assume the result holds for solvable groups of length $n-1$ and let
% 	$G$ be a solvable group with a sequence~\eqref{eq:resoluble}. Since $G_{n-1}$ is 
% 	a solvable torsion group, the inductive hypothesis implies that 
% 	$G_{n-1}$ is locally finite. Since $G/G_{n-1}$ is an abelian torsion group, 
% 	it is locally finite. The result now follows from Proposition \ref{pro:exact_LI}.
% \end{proof}

\begin{bonus}
    \label{bonus:solvable=>LF}
    Prove that solvable torsion groups are locally finite.  
\end{bonus}

% \begin{sol}{bonus:solvable=>LF}
% 	We proceed by induction on $n$, the length of the filtration. 
% 	If $n=1$, then $G$ is finite because it is abelian and a torsion group.
% 	Now assume the result holds for solvable groups of length $n-1$ and let
% 	$G$ be a solvable group with a filtration of length $n$. Since $G_{n-1}$ is 
% 	a solvable torsion group, the inductive hypothesis implies that 
% 	$G_{n-1}$ is locally finite. Since $G/G_{n-1}$ is an abelian torsion group, 
% 	it is locally finite. The result now follows from Proposition \ref{pro:exact_LI}.
% \end{proof}


\begin{theorem}[Herstein]
\index{Herstein's theorem}
	If $G$ is a locally finite group, then $K[G]$ is algebraic. Conversely, if 
	$K[G]$ is algebraic and $K$ has characteristic zero, then $G$ 
	is locally finite. 
\end{theorem}

\begin{proof}
	Assume that $G$ is locally finite. Let $\alpha\in K[G]$. The subgroup 
	$H=\langle\supp\alpha\rangle$ is finite, as it is finitely generated. Since 
	$\alpha\in K[H]$ and $\dim_KK[H]<\infty$, the set 
	$\{1,\alpha,\alpha^2,\dots\}$ is linearly dependent. Thus $\alpha$ is
	algebraic over $K$.

	Let $\{x_1,\dots,x_m\}$ be a finite subset of $G$. Adding inverses if needed,
	we may assume that $\{x_1,\dots,x_m\}$ generates the subgroup 
	$H=\langle x_1,\dots,x_m\rangle$ as a semigroup. Let 
	\[
	\alpha=x_1+\dots+x_m\in K[G].
	\]
	Since $\alpha$ is algebraic over $K$, 
	there exist $b_0,b_1,\dots,b_{n+1}\in K$ such that 
	\[
	b_0+b_1\alpha+\cdots+b_{n+1}\alpha^{n+1}=0,
	\]
	where $b_{n+1}\ne 0$. We can rewrite this as 
	\[
		\alpha^{n+1}=a_0+a_1\alpha+\cdots+a_n\alpha^n
	\]
	for some $a_0,\dots,a_n\in K$. 
	Note that 
	\[
		\alpha^{k}=(x_1+\cdots+x_m)^{k}
		=\sum x_{i_1}\cdots x_{i_{k}}
	\]
	for all $k$. 
	Two words $x_{i_1}\cdots x_{i_{k}}$ and 
	$x_{j_1}\cdots x_{j_{l}}$ could represent the same element of the group $H$. In this case, 
	the coefficient of $x_{i_1}\cdots x_{i_{k}}=x_{j_1}\cdots x_{j_{l}}$ 
	in $\alpha^{k}$ will be a positive integer $\geq2$. 
 
        Let $w=x_{i_1}\cdots
	x_{i_{n+1}}\in H$ be a word of length $n+1$. 
		Since $K$
	is of characteristic zero, it follows that $w\in\supp(\alpha^{n+1})$. Since, moreover,  
	$\alpha^{n+1}=\sum_{j=0}^na_j\alpha^j$, it follows that 
	$w\in\supp(\alpha^j)$ for some $j\in\{0,\dots,n\}$. Thus each
	word in the letters $x_j$ of length $n+1$ can be written as a word in the letters $x_j$ of 
	length $\leq n$. Therefore $H$ is finite and hence $G$ is locally finite. 
\end{proof}


\section{Formanek's theorem, I}

We start with some exercises. 

\begin{exercise}
\label{xca:invertible_algebraic}
	Let $K$ be a field. Let $A$ be a $K$-algebra algebraic over $K$ and $a\in A$.
	\begin{enumerate}
		\item $a$ is a left zero divisor if and only if $a$ is a right zero divisor.
		\item $a$ is left invertible if and only if $a$ is right invertible.
		\item $a$ is invertible if and only if $a$ is not a zero divisor.
	\end{enumerate}
\end{exercise}

\begin{exercise}
	\label{exa:norma}
	For $\alpha=\sum_{g\in G}\alpha_gg\in\C[G]$ let $|\alpha|=\sum_{g\in
	G}|\alpha_g|\in\R$. Prove the following statements:
	\begin{enumerate}
		%\item $|\trace(\alpha)|\leq |\alpha|$, 
		\item $|\alpha+\beta|\leq|\alpha|+|\beta|$, and 
		\item $|\alpha\beta|\leq|\alpha||\beta|$ 
	\end{enumerate}
	for all $\alpha,\beta\in\C[G]$.
\end{exercise}

\begin{theorem}[Formanek]
	\label{thm:FormanekQ}
	\index{Formanek's theorem}
	Let $G$ be a group. If every element of $\Q[G]$ is invertible or 
	a zero divisor, then $G$ is locally finite. 
\end{theorem}

\begin{proof}
	Let $\{x_1,\dots,x_n\}$ be a finite subset of $G$. Adding inverses if needed, we may assume that 
	$\{x_1,\dots,x_n\}$ generates the subgroup
	$H=\langle x_1,\dots,x_n\rangle$ as a semigroup. Let 
	\[
		\alpha=\frac{1}{2n}(x_1+\cdots+x_n)\in\Q[G]
	\]
    Note that $|\alpha|\leq 1/2$. 
	We claim that $1-\alpha\in\Q[G]$ is invertible. If not, then it is a zero divisor. If there exists 
	$\delta\in\Q[G]$ such that $\delta(1-\alpha)=0$, then 
	$\delta=\delta\alpha$. Since  
	\[
		|\delta|=|\delta\alpha|\leq|\delta||\alpha|\leq|\delta|/2,
	\]
	it follows that $\delta=0$. Similarly, $(1-\alpha)\delta=0$ implies
	$\delta=0$. 
	
	Let $\beta=(1-\alpha)^{-1}\in\Q[G]$.  For each $k$ let  
	\[
		\gamma_k=(1+\alpha+\cdots+\alpha^k)-\beta.
	\]
	Then 
	\begin{align*}
		\gamma_k(1-\alpha)&=(1+\alpha+\cdots+\alpha^k-\beta)(1-\alpha)\\
		&=(1+\alpha+\cdots+\alpha^k)(1-\alpha)-\beta(1-\alpha)=-\alpha^{k+1}
	\end{align*}
	and thus  
	$\gamma_k=-\alpha^{k+1}\beta$. Since  
	\[
		|\gamma_k|=|-\alpha^{k+1}\beta|\leq|\beta||\alpha^{k+1}|\leq\frac{|\beta|}{2^{k+1}},
	\]
	it follows that $\lim_{k\to\infty}|\gamma_k|=0$. 

	We now prove that $H\subseteq\supp\beta$. This will finish the proof of the theorem, 
	as $\supp\beta$ is a finite subset of $G$ by definition. If
	$H\not\subseteq\supp\beta$, let $h\in H\setminus\supp\beta$.  Assume that 
    $h=x_{i_1}\cdots x_{i_m}$ is a word in the letters $x_j$ of length $m$. Let 
    $c_j$ be the coefficient of $h$ in $\alpha^j$. Then $c_0+\cdots+c_k$ is the 
	coefficient of $h$ in $\gamma_k$, but 
	\[
		|\gamma_k|\geq c_0+c_1+\cdots+c_k\geq c_m>0
	\]
	for all $k\geq m$, as each $c_j$ is non-negative, a contradiction to 
	$|\gamma_k|\to 0$ if $k\to\infty$.
\end{proof}



\section{Tensor products}

The \emph{tensor product} of the vector spaces (over $K$) $U$ and $V$ 
is the quotient vector space $K[U\times V]/T$, where $K[U\times V]$ 
is the vector space with basis 
\[
\{(u,v):u\in U,v\in V\}
\]
and $T$ is the subspace 
generated by elements of the form 
\[
		(\lambda u+\mu u',v)-\lambda(u,v)-\mu(u',v),\quad
		(u,\lambda v+\mu v')-\lambda(u,v)-\mu(u,v')
	\]
for $\lambda,\mu\in K$, $u,u'\in U$ and $v,v'\in V$.
The tensor product of $U$ and $V$ will be denoted by $U\otimes_KV$ or 
$U\otimes V$ when the base field is clear from the context. For $u\in U$ and 
$v\in V$ we write $u\otimes v$ to denote the coset $(u,v)+T$.

\begin{theorem}
	Let $U$ and $V$ be vector spaces. Then there exists a bilinear map 
	\[
    U\times V\to U\otimes V,\quad (u,v)\mapsto u\otimes v,
    \]
    such that 
	each element of $U\otimes V$ is a finite sum of the form 
	\[
		\sum_{i=1}^N u_i\otimes v_i
	\]
	for some $u_1,\dots,u_N\in U$ and $v_1,\dots,v_N\in V$. 
	Moreover, if $W$ is a vector space and 
 \[
 \beta\colon U\times V\to W
 \]
 is a bilinear map, 
	there exists a linear map 
	$\overline{\beta}\colon U\otimes V\to W$ such that $\overline{\beta}(u\otimes
	v)=\beta(u,v)$ for all $u\in U$ and $v\in V$.
\end{theorem}

\begin{proof}
    By definition, the map
    \[
	U\times V\to U\otimes V,\quad
	(u,v)\mapsto u\otimes v,
	\]
	is bilinear. From the definitions, it follows that
	$U\otimes V$ is a finite linear combination of elements of the form 
	$u\otimes v$, where $u\in U$ and $v\in V$. Since $\lambda(u\otimes
	v)=(\lambda u)\otimes v$ for all $\lambda\in K$, the first claim follows.

	Since the elements of $U\times V$ form a basis of $K[U\times V]$, there exists
	a linear map 
	\[
		\gamma\colon K[U\times V]\to W,\quad
	\gamma(u,v)=\beta(u,v). 
	\]
	Since $\beta$ is bilinear by assumption, $T\subseteq\ker\gamma$. It follows that there exists 
	a linear map $\overline{\beta}\colon U\otimes V\to
	W$ such that  
	\[
	\begin{tikzcd}
		K[U\times V] \arrow[r]\arrow[d] & W \\
		U\otimes V\arrow[ur, dashrightarrow]
	\end{tikzcd}
	\]
	commutes. In particular, $\overline{\beta}(u\otimes v)=\beta(u,v)$. 
\end{proof}

\begin{exercise}
	\label{xca:tensorial_unicidad}
	Prove that the properties of the previous theorem characterize tensor products up to isomorphism. 
\end{exercise}

Some properties:
%Observemos
%que todo elemento de $U\otimes V$ es una suma finita
%de la forma 
%\[
%	\sum_{i=1}^N u_i\otimes v_i
%\]
%para $N\in\N$, $u_i\in U$ y $v_i\in V$. Esta expresión no es única. Vale además
%que $u\otimes 0=0=0\otimes v$ para todo $u\in U$ y $v\in V$.

\begin{proposition}
	Let $\varphi\colon U\to U_1$ and $\psi\colon V\to V_1$ be linear maps. There
	exists a unique linear map 
	$\varphi\otimes\psi\colon U\otimes V\to U_1\otimes V_1$ such that
	\[
		(\varphi\otimes\psi)(u\otimes v)=\varphi(u)\otimes\psi(v)
	\]
	for all $u\in U$ and $v\in V$.
\end{proposition}

\begin{proof}
	Since $U\times V\to U_1\otimes V_1$,
	$(u,v)\mapsto\varphi(u)\otimes\psi(v)$, is bilinear, there exists a linear map
	$U\otimes V\to U_1\otimes V_1$, $u\otimes
	v\to\varphi(u)\otimes\psi(v)$. Thus 
	\[
		\sum u_i\otimes v_i\mapsto\sum\varphi(u_i)\otimes\psi(v_i)
	\]
	is well-defined. 
\end{proof}

\begin{exercise}
    Prove the following statements:
	\begin{enumerate}
		\item $(\varphi\otimes\psi)(\varphi'\otimes\psi')=(\varphi\varphi')\otimes(\psi\psi')$.
		\item If $\varphi$ and $\psi$ are isomorphisms, then 
			$\varphi\otimes\psi$ is an isomorphism. 
		\item $(\lambda\varphi+\lambda'\varphi')\otimes\psi=\lambda\varphi\otimes\psi+\lambda'\varphi'\otimes\psi$.
		\item $\varphi\otimes(\lambda\psi+\lambda'\psi')=\lambda\varphi\otimes\psi+\lambda'\varphi\otimes\psi'$.
		\item If $U\simeq U_1$ and $V\simeq V_1$, then $U\otimes V\simeq U_1\otimes V_1$.
	\end{enumerate}
\end{exercise}

The following proposition is extremely useful:

\begin{proposition}
	If $U$ and $V$ are vector spaces, then  
	$U\otimes V\simeq V\otimes U$.
\end{proposition}

\begin{proof}
	Since $U\times V\to V\otimes U$, $(u,v)\mapsto v\otimes u$, is bilinear, there exists 
	a linear map 
    \[
    U\otimes V\to V\otimes U,\quad u\otimes
	v\mapsto v\otimes u.
    \]
    Similarly, there exists a linear map 
	\[
    V\otimes U\to U\otimes V,\quad  v\otimes u\mapsto
	u\otimes v.
    \]
    Thus $U\otimes V\simeq V\otimes U$.
\end{proof}

\begin{exercise}
	\label{xca:UxVxW}
    Prove that $(U\otimes V)\otimes W\simeq U\otimes(V\otimes W)$.
\end{exercise}

\begin{exercise}
	\label{xca:UxK}
	Prove that $U\otimes K\simeq U\simeq K\otimes U$.
\end{exercise}

\begin{proposition}
	\label{pro:U_LI}
	Let $U$ and $V$ be vector spaces. 
	If $\{u_1,\dots,u_n\}$ is a linearly independent subset of $U$ and 
	$v_1,\dots,v_n\in V$ is such that $\sum_{i=1}^n u_i\otimes v_i=0$, then 
    $v_i=0$ for all $i\in\{1,\dots,n\}$.
\end{proposition}

\begin{proof}
	Let $i\in\{1,\dots,n\}$ and 
	\[
	f_i\colon U\to K,
	\quad
	f_i(u_j)=\delta_{ij}=\begin{cases}
	1 & \text{if $i=j$},\\
	0 & \text{otherwise}.
	\end{cases}
	\]
	Since the map 
 \[
 U\times V\to V,\quad (u,v)\mapsto f_i(u)v,
 \]
 is bilinear, there exists 
	a linear map 
	$\alpha_i\colon U\otimes V\to V$ such that $\alpha_i(u\otimes
	v)=f_i(u)v$. Thus
	\[
		v_i=\sum_{j=1}^n\alpha_i(u_j\otimes v_j)=\alpha_i\left(\sum_{j=1}^nu_j\otimes v_j\right)=0.\qedhere
	\]
\end{proof}

\begin{exercise}
	\label{xca:uxv=0}
	Prove that $u\otimes v=0$ and $v\ne 0$ imply $u=0$.
\end{exercise}

\begin{theorem}
    Let $U$ and $V$ be vector spaces. 
	If $\{u_i:i\in I\}$ is a basis of $U$ and $\{v_j:j\in J\}$ is a basis of $V$, then 
	$\{u_i\otimes v_j:i\in I,j\in J\}$ is a basis of $U\otimes
	V$.
\end{theorem}

\begin{proof}
	The $u_i\otimes v_j$ are generators of $U\otimes V$, as  
    $u=\sum_i\lambda_iu_i$ and $v=\sum_j\mu_jv_j$ imply 
	\[
    u\otimes v=\sum_{i,j}\lambda_i\mu_ju_i\otimes v_j.
    \]
	We now prove that the $u_i\otimes v_j$ are linearly independent. We need to show that
	each finite subset of the $u_i\otimes v_j$
	is linearly independent. If $\sum_k\sum_l\lambda_{kl}u_{i_k}\otimes
	v_{j_l}=0$, then 
	\[
0=\sum_{k}u_{i_k}\otimes\left(\sum_{l}\lambda_{kl}v_{j_l}\right).\]
    Since  
	the $u_{i_k}$ are linearly independent, Proposition~\ref{pro:U_LI}
	implies that $\sum_{l}\lambda_{kl}v_{j_l}=0$. Thus $\lambda_{kl}=0$ for all 
	$k,l$, as the $v_{j_l}$ are linearly independent. 
\end{proof}

If $U$ and $V$ are finite-dimensional vector spaces, then 
\[
	\dim(U\otimes V)=(\dim U)(\dim V).
\]

\begin{corollary}
	If $\{u_i:i\in I\}$ is a basis of $U$, then every element of $U\otimes V$
	can be written uniquely as a finite sum $\sum_{i}u_i\otimes v_i$.
\end{corollary}

\begin{proof}
	Every element of the tensor product $U\otimes V$ is a finite sum 
	of the form $\sum_i x_i\otimes y_i$, where $x_i\in U$ and $y_i\in V$. If  
	$x_i=\sum_j\lambda_{ij}u_j$, then 
	\[
		\sum_i x_i\otimes y_i=\sum_i\left(\sum_j\lambda_{ij}u_j\right)\otimes y_i		
		=\sum_j u_j\otimes\left(\sum_i\lambda_{ij}y_i\right).\qedhere 
	\]
\end{proof}

%\begin{corollary}
%	Todo elemento no nulo de $U\otimes V$ puede escribirse como una suma finita
%	$\sum_{i=1}^N u_i\otimes v_i$ para un conjuntos $\{u_i:1\leq i\leq
%	N\}\subseteq U$ y $\{v_i:1\leq i\leq N\}\subseteq V$ linealmente
%	independientes.
%\end{corollary}
%
%\begin{proof}
%	tomar $N$ minimal	
%\end{proof}

\begin{exercise}
\label{xca:tensor_algebras}
    Let $A$ and $B$ be algebras. Prove that $A\otimes B$ 
    is an algebra with 
	\[
		(a\otimes b)(x\otimes y)=ax\otimes by.
	\]
\end{exercise}

% \begin{proof}
% 	Para $x\in A$, $y\in B$ consideramos $R_x\otimes R_y\in\End_K(A\otimes B)$.
% 	Como la función $A\times B\to\End_K(A\otimes B)$, $(x,y)\mapsto R_x\otimes
% 	R_y$, es bilineal, existe una función lineal $\varphi\colon A\otimes
% 	B\to\End_K(A\otimes B)$, $\varphi(x\otimes y)=R_x\otimes R_y$. Para $u,v\in A\otimes B$ definimos
% 	\[
% 		uv=\varphi(v)(u).
% 	\]
% 	Esta operación es bilineal pues por ejemplo
% 	\[
% 		u(v+w)=\varphi(v+w)(u)=(\varphi(v)+\varphi(w))(u)=\varphi(v)(u)+\varphi(w)(u)=uv+uw.
% 	\]
% 	Además
% 	$(a\otimes b)(x\otimes y)=\varphi(x\otimes y)(a\otimes b)=(R_x\otimes R_y)(a\otimes b)=ax\otimes by$.
% 	Un cálculo sencillo muestra que este producto es asociativo.
% \end{proof}

\begin{exercise}
    Let $K$ be a field and $A,B,C$ be $K$-algebras. 
    Prove the following statements:
	\begin{enumerate}
		\item $A\otimes B\simeq B\otimes A$.
		\item $(A\otimes B)\otimes C\simeq A\otimes(B\otimes C)$.
		\item $A\otimes K\simeq A\simeq K\otimes A$.
		\item If $A\simeq A_1$ and $B\simeq B_1$, then $A\otimes B\simeq A_1\otimes B_1$.
	\end{enumerate}
\end{exercise}

Some examples:

\begin{proposition}
	If $G$ and $H$ are groups, then $K[G]\otimes K[H]\simeq K[G\times H]$.
\end{proposition}

\begin{proof}
	The set $\{g\otimes h:g\in G,h\in H\}$ is a basis of $K[G]\otimes K[H]$ and 
	the elements of $G\times H$ form a basis of $K[G\times H]$. There exists a linear isomorphism 
	\[
	K[G]\otimes K[H]\to K[G\times H], 
	\quad 
	g\otimes h\mapsto (g,h),
	\]
	that is multiplicative. Thus $K[G]\otimes K[H]\simeq K[G\times H]$
	as algebras. 
\end{proof}

\begin{proposition}
\label{pro:AKX=AX}
	If $A$ is an algebra, then $A\otimes K[X]\simeq A[X]$.	
\end{proposition}

\begin{proof}
	Each element of the tensor product $A\otimes K[X]$ can be written uniquely as a finite sum of
	the form $\sum a_i\otimes X^i$. Routine calculations show that 
	$A\otimes K[X]\mapsto A[X]$, $\sum a_i\otimes X^i\mapsto \sum a_iX^i$, is a 
	linear algebra isomorphism. 
\end{proof}

\begin{exercise}
\label{xca:AM=MA}
	Prove that if $A$ is an algebra, then $A\otimes M_n(K)\simeq M_n(A)$. In
	particular, $M_n(K)\otimes M_m(K)\simeq M_{nm}(K)$.
\end{exercise}

Proposition \ref{pro:AKX=AX} and Exercise \ref{xca:AM=MA} 
are examples of a procedure known as \emph{scalar extensions}. 

\begin{theorem}
\label{thm:extensions_scalars}
	Let $A$ be an algebra over $K$ and $E$ be an extension of $K$ (this just simply means that
	$K$ is a subfield of $E$). Then 
	$A^E=E\otimes_KA$ is an algebra over $E$ with respect to
	the scalar multiplication 
	\[
		\lambda(\mu\otimes a)=(\lambda\mu)\otimes a,
	\]
	for all $\lambda,\mu\in E$ and $a\in A$.
\end{theorem}

\begin{proof}
	Let $\lambda\in E$. Since $E\times A\to E\otimes_KA$,
	$(\mu,a)\mapsto (\lambda\mu)\otimes a$, is $K$-bilinear, there exists 
	a linear map $E\otimes_KA\to E\otimes_KA$, $\mu\otimes a\mapsto
	(\lambda\mu)\otimes a$. The scalar multiplication is then well-defined and 
	\[
	\lambda(u+v)=\lambda u+\lambda v
	\]
	for all $\lambda\in E$ and $u,v\in E\otimes_KA$. Moreover, 
	\[
	(\lambda+\mu)u=\lambda u+\mu u,
	\quad
	(\lambda\mu)u=\lambda(\mu u),
	\quad
	\lambda(uv)=(\lambda u)v=u(\lambda v)
	\]
	for all $u,v\in E\otimes_KA$ and $\lambda,\mu\in E$.
\end{proof}

\begin{exercise}
    Prove the following statements:
    \begin{enumerate}
		\item $\{1\}\otimes A$ is a subalgebra of $A^E$ isomorphic to $A$.
		\item If $\{a_i:i\in I\}$ is a basis of $A$, then $\{1\otimes a_i:i\in
			I\}$ is a basis of $A^E$.
	\end{enumerate}
\end{exercise}

\begin{exercise}
	Prove that if $G$ is a group and $K$ is a subfield of $E$, then
	\[
	E\otimes_K K[G]\simeq E[G].
	\]
\end{exercise}

\section{Formanek's theorem, II}

The combination of technique known as extensions of scalars we have seen in the previous section
and Formanek's theorem for rational group algebras 
yield the following general result. 

\begin{theorem}[Formanek]
	\index{Formanek's theorem}
	Let $K$ be a field of characteristic zero and let $G$ be a group. 
	If every element of $K[G]$ is invertible or a zero divisor, 
	then $G$ is locally finite. 
\end{theorem}

\begin{proof}
	Since $K$ is of characteristic zero, $\Q\subseteq K$. Then $K[G]\simeq
	K\otimes_{\Q}\Q[G]$. Note that each $\beta\in K\otimes_{\Q}\Q[Q]$ can be written
	uniquely as 
	\[
		\beta=1\otimes\beta_0+\sum k_i\otimes\beta_i,
	\]
	where $\{1,k_1,k_2,\dots,\}$ is a basis of $K$ as a $\Q$-vector space. 
	Let $\alpha\in\Q[G]\subseteq K[G]$. By assumption, every element of $K[G]$ is 
    invertible or a zero divisor. Assume first that $\alpha$ is invertible. 
    Let $\beta\in K[G]\simeq K\otimes_{\Q}\Q[Q]$ be such that $\alpha\beta=1$. Since
	\[
	1\otimes 1=(1\otimes\alpha)\beta=1\otimes \alpha\beta_0+\sum k_i\otimes \alpha\beta_i,
	\]
	it follows that $\alpha\beta_0=1$. Thus $\alpha$ is invertible in $\Q[G]$. Similarly, if
    $\alpha$ is a zero divisor, then 
	$\alpha\beta=0$ and $\alpha\beta_j=0$ for all $j$. Thus $\alpha$ is a zero divisor in $\Q[G]$. 
    Hence 
	each $\alpha\in\Q[G]$ is invertible or a zero divisor, so Formanek's theorem 
	for $\Q$ applies.
\end{proof}


